% GENERATED FILE. Edit canonical lessons, then run npm run generate:books.
% canonical-source-hash: fnv1a64:86b4c8715ab682e2
% canonical-lessons: UR-C297-kaimp, UR-C297-rahat, UR-C297-khaima, UR-C297-panah, UR-C297-mutasir

\chapter{The Relief Camp}
\label{ch:ur-the-relief-camp}
\hlchaptermodality{297}

\begin{chapteropening}
\textbf{By the end of this chapter:} \emph{I can describe life in a relief camp.}

Complete the last lesson of chapter 297: I can describe life in a relief camp.
\end{chapteropening}

\section[\texorpdfstring{\ur{کیمپ} (kaimp)}{کیمپ (kaimp)}]{\ur{کیمپ} (kaimp) — a relief camp; a camp}

\label{lesson:UR-C297-kaimp}



\begin{quote}
Before the new one: say the Urdu for to dart, then the Urdu for safe and sound.
\end{quote}

\subsection*{You'll want to know: \ur{کیمپ}}
\textbf{\ur{کیمپ}} — \emph{kaimp} — \textquotedblleft{}a relief camp; a camp\textquotedblright{}.

\begin{quote}
\ur{شہر} \ur{کے} \ur{باہر} \ur{ایک} \ur{کیمپ} \ur{تھا۔}
\end{quote}

\begin{quote}
\emph{shahr ke bāhar ek kaimp thā.}
\end{quote}

\begin{quote}
There was a camp outside the city.
\end{quote}

Masculine.

\begin{cousinweb}
is borrowed from English \emph{camp}. After a flood or an earthquake, people live in a \textbf{\ur{کیمپ}} for a time.
\end{cousinweb}

\begin{grammarlens}[title={What you've built}]
One more word for telling what happened.
\end{grammarlens}

\subsection*{Your turn}
\begin{itemize}
\raggedright
  \item \emph{Say it:} \emph{kaimp}
  \item \emph{Say it:} \emph{kaimp}, once more
  \item \emph{Say it:} \emph{shahr ke bāhar ek kaimp thā.} — There was a camp outside the city.
  \item \emph{Recall:} say the Urdu for to dart, then the Urdu for safe and sound, then say \emph{kaimp} again
\end{itemize}

\begin{tcolorbox}[breakable,colback=teal!4,colframe=teal!35!black,title={Before you move on}]
What does \ur{کیمپ} mean? (\textquotedblleft{}A relief camp; a camp\textquotedblright{}.) Say it once more.
\end{tcolorbox}
\section[\texorpdfstring{\ur{راحت} (rāhat)}{راحت (rāhat)}]{\ur{راحت} (rāhat) — relief}

\label{lesson:UR-C297-rahat}



\begin{quote}
Before the new one: say the Urdu for safe and sound, then the Urdu for a relief camp.
\end{quote}

\subsection*{You'll want to know: \ur{راحت}}
\textbf{\ur{راحت}} — \emph{rāhat} — \textquotedblleft{}relief\textquotedblright{}.

\begin{quote}
\ur{کیمپ} \ur{میں} \ur{لوگوں} \ur{کو} \ur{راحت} \ur{ملی۔}
\end{quote}

\begin{quote}
\emph{kaimp meṅ logoṅ ko rāhat milī.}
\end{quote}

\begin{quote}
People found relief in the camp.
\end{quote}

Feminine.

\begin{cousinweb}
\textbf{\ur{راحت}} (\emph{rāhat}) is usually traced, through Persian, to Arabic \emph{rāḥa}, \textquotedblleft{}rest, ease\textquotedblright{}. It is feminine: \textbf{\ur{راحت} \ur{ملی}}.
\end{cousinweb}

\begin{grammarlens}[title={What you've built}]
One more word for telling what happened.
\end{grammarlens}

\subsection*{Your turn}
\begin{itemize}
\raggedright
  \item \emph{Say it:} \emph{rāhat}
  \item \emph{Say it:} \emph{rāhat}, once more
  \item \emph{Say it:} \emph{kaimp meṅ logoṅ ko rāhat milī.} — People found relief in the camp.
  \item \emph{Recall:} say the Urdu for safe and sound, then the Urdu for a relief camp, then say \emph{rāhat} again
\end{itemize}

\begin{tcolorbox}[breakable,colback=teal!4,colframe=teal!35!black,title={Before you move on}]
What does \ur{راحت} mean? (\textquotedblleft{}Relief\textquotedblright{}.) Say it once more.
\end{tcolorbox}
\section[\texorpdfstring{\ur{خیمہ} (khaimā)}{خیمہ (khaimā)}]{\ur{خیمہ} (khaimā) — a tent}

\label{lesson:UR-C297-khaima}



\begin{quote}
Before the new one: say the Urdu for a relief camp, then the Urdu for relief.
\end{quote}

\subsection*{You'll want to know: \ur{خیمہ}}
\textbf{\ur{خیمہ}} — \emph{khaimā} — \textquotedblleft{}a tent\textquotedblright{}.

\begin{quote}
\ur{ہر} \ur{خاندان} \ur{کو} \ur{ایک} \ur{خیمہ} \ur{ملا۔}
\end{quote}

\begin{quote}
\emph{har khāndān ko ek khaimā milā.}
\end{quote}

\begin{quote}
Each family got a tent.
\end{quote}

Masculine.

\begin{cousinweb}
\textbf{\ur{خیمہ}} (\emph{khaimā}) is usually traced, through Persian, to Arabic \emph{khayma}, \textquotedblleft{}a tent\textquotedblright{}.
\end{cousinweb}

\begin{grammarlens}[title={What you've built}]
One more word for telling what happened.
\end{grammarlens}

\subsection*{Your turn}
\begin{itemize}
\raggedright
  \item \emph{Say it:} \emph{khaimā}
  \item \emph{Say it:} \emph{khaimā}, once more
  \item \emph{Say it:} \emph{har khāndān ko ek khaimā milā.} — Each family got a tent.
  \item \emph{Recall:} say the Urdu for a relief camp, then the Urdu for relief, then say \emph{khaimā} again
\end{itemize}

\begin{tcolorbox}[breakable,colback=teal!4,colframe=teal!35!black,title={Before you move on}]
What does \ur{خیمہ} mean? (\textquotedblleft{}A tent\textquotedblright{}.) Say it once more.
\end{tcolorbox}
\section[\texorpdfstring{\ur{پناہ} (panāh)}{پناہ (panāh)}]{\ur{پناہ} (panāh) — refuge; shelter}

\label{lesson:UR-C297-panah}



\begin{quote}
Before the new one: say the Urdu for relief, then the Urdu for a tent.
\end{quote}

\subsection*{You'll want to know: \ur{پناہ}}
\textbf{\ur{پناہ}} — \emph{panāh} — \textquotedblleft{}refuge; shelter\textquotedblright{}.

\begin{quote}
\ur{لوگوں} \ur{نے} \ur{اسکول} \ur{میں} \ur{پناہ} \ur{لی۔}
\end{quote}

\begin{quote}
\emph{logoṅ ne iskūl meṅ panāh lī.}
\end{quote}

\begin{quote}
People took shelter in the school.
\end{quote}

Feminine.

\begin{cousinweb}
\textbf{\ur{پناہ}} (\emph{panāh}) is usually traced to Persian \emph{panāh}, \textquotedblleft{}shelter, protection\textquotedblright{}. \textbf{\ur{پناہ} \ur{لی}} is took refuge.
\end{cousinweb}

\begin{grammarlens}[title={What you've built}]
One more word for telling what happened.
\end{grammarlens}

\subsection*{Your turn}
\begin{itemize}
\raggedright
  \item \emph{Say it:} \emph{panāh}
  \item \emph{Say it:} \emph{panāh}, once more
  \item \emph{Say it:} \emph{logoṅ ne iskūl meṅ panāh lī.} — People took shelter in the school.
  \item \emph{Recall:} say the Urdu for relief, then the Urdu for a tent, then say \emph{panāh} again
\end{itemize}

\begin{tcolorbox}[breakable,colback=teal!4,colframe=teal!35!black,title={Before you move on}]
What does \ur{پناہ} mean? (\textquotedblleft{}Refuge; shelter\textquotedblright{}.) Say it once more.
\end{tcolorbox}
\section[\texorpdfstring{\ur{متاثر} (mutāsir)}{متاثر (mutāsir)}]{\ur{متاثر} (mutāsir) — affected (by a disaster)}

\label{lesson:UR-C297-mutasir}



\begin{quote}
Before the new one: say the Urdu for a tent, then the Urdu for refuge.
\end{quote}

\subsection*{You'll want to know: \ur{متاثر}}
\textbf{\ur{متاثر}} — \emph{mutāsir} — \textquotedblleft{}affected (by a disaster)\textquotedblright{}.

\begin{quote}
\ur{سیلاب} \ur{سے} \ur{دو} \ur{گاؤں} \ur{متاثر} \ur{ہوئے۔}
\end{quote}

\begin{quote}
\emph{sailāb se do gāoṅ mutāsir hue.}
\end{quote}

\begin{quote}
Two villages were affected by the flood.
\end{quote}

\begin{cousinweb}
\textbf{\ur{متاثر}} (\emph{mutāsir}) is usually traced, through Persian, to Arabic \emph{mutaʾaththir}, \textquotedblleft{}one that has taken an effect\textquotedblright{}, a word from the same family as \textbf{\ur{اثر}}.
\end{cousinweb}

\begin{grammarlens}[title={What you've built}]
One more word for telling what happened.
\end{grammarlens}

\subsection*{Your turn}
\begin{itemize}
\raggedright
  \item \emph{Say it:} \emph{mutāsir}
  \item \emph{Say it:} \emph{mutāsir}, once more
  \item \emph{Say it:} \emph{sailāb se do gāoṅ mutāsir hue.} — Two villages were affected by the flood.
  \item \emph{Recall:} say the Urdu for a tent, then the Urdu for refuge, then say \emph{mutāsir} again
\end{itemize}

\begin{tcolorbox}[breakable,colback=teal!4,colframe=teal!35!black,title={Before you move on}]
What does \ur{متاثر} mean? (\textquotedblleft{}Affected (by a disaster)\textquotedblright{}.) Say it once more.
\end{tcolorbox}
